\documentclass[12pt]{article}

\usepackage{graphicx}
\usepackage{titling}
\usepackage{amsmath}
\usepackage{amssymb}
% \usepackage{tensor}
% \usepackage{ragged2e}
\usepackage{lastpage}
\usepackage[margin=2.5cm]{geometry}
\usepackage{blindtext}
\usepackage{titlesec}
\usepackage{amssymb}
\usepackage[colorlinks=true, linkcolor=blue]{hyperref}
% \usepackage{changepage}
\begin{document}


\begin{center}
    {\LARGE Stochastic Barrel Pyrolysis \\ [0.5cm]
    \large Joe Salter}
\end{center}


\section{Introduction}
We consider PuO and another material which is combustible. Considerations at the start will be general, but we will ultimately impose a slab geometry. The equations to be coupled are the sourced heat conduction equation and the time-dependent neutron diffusion equation with delayed precursors.

\subsection{Heat Conduction}
The sourced conduction equation takes the form 
\begin{equation}
    \rho (\mathbf{r}) C_p (\mathbf{r}) \frac{\partial T(\bold{r},t)}{\partial t}=\bold{\nabla} \cdot (K(\mathbf{r})\bold{\nabla} T) + S(\bold{r},t). \label{SCE}
\end{equation}
We split the source into two parts, 
\begin{equation}
    S(\bold{r},t)=S_{fission}+S_{combustion}.
\end{equation}
$S_{combustion}$ arises from decomposition, and is given by 
\begin{equation}
    S_{combustion}(\bold{r},t)=q(\mathbf{r})\left( -\frac{\partial \omega(\bold{r},t)}{\partial t} \right)
\end{equation}
with 
\begin{equation}
    -\frac{\partial \omega(\bold{r},t)}{\partial t}=k(\mathbf{r})\omega (\mathbf{r},t) exp\left(-\frac{E_{act} (\mathbf{r})}{R_g T(\mathbf{r},t)} \right).
\end{equation}
The fission source term has units of $Jm^{-3}s^{-1}$, so takes the form $E_f \times n \times v \times \Sigma_f$. Hence 
\begin{equation}
    S_{fission}=E_f \int_0 ^\infty \Sigma_f(E)\phi(\mathbf{r}, E, t)dE.
\end{equation}

\subsection{Neutron diffusion}
The time dependent NDE takes the form 
\begin{align}
    \frac{1}{v(E)} \frac{\partial \phi (\bold{r},E,t)}{\partial t}&=\bold{\nabla} \cdot (D(\bold{r}, E)\bold{\nabla} \phi)-\Sigma_a (\bold{r},E)\phi+\int_{0}^{\infty}\Sigma_s(\bold{r}, E'\rightarrow E)\phi(\bold{r},E',t)dE' \notag \\ & +\chi_p(E) \int_0 ^{\infty} \nu(E') (1-\beta)\Sigma_f (\bold{r},E')\phi(\bold{r},E',t)dE' +\sum_{i=1}^{I} \chi_{d,i}(E)\lambda_i C_i(\bold{r},t) \label{TDNDE}
\end{align}
with the precursor equations 
\begin{equation}
    \frac{\partial C_i({\bold{r},t)}}{\partial t}=\beta_i \int_0^\infty \nu(E) \Sigma_f (\bold{r},E) \phi(\bold{r}, E,t)dE-\lambda_i C_i .
\end{equation}

\section{Solving}
When solving the coupled thermal and neutronic equations, we should for each timestep first solve the thermal equations, then update the necessary cross sections and coefficients, and solve the neutronic equation. For this reason, a dependence on $t$ may in some cases be seen as a dependence on $T$.

\section{A Stochastic Ansatz}
\subsection{The Formalism}
Let us now introduce the machinery to deal with randomly dispersed mixtures. We define the random function $F(r,\vartheta)$ as 
\begin{equation}
    F(r,\vartheta)=1 \hspace{0.25cm}\text{If $r$ is in phase 1,} \hspace{0.25cm} F(r, \vartheta)=0 \hspace{0.25cm} \text{If $r$ is in phase 2}
\end{equation}
where $\vartheta$ is a set of realizations. We can write a random function as 
\begin{equation}
    \Psi(r, \vartheta)=\Psi_1 F+ \Psi_2 (1-F).
\end{equation}
We identify the statistical average of $F$ as the volume fraction of phase 1, and vice versa for $1-F$ :
\begin{align}
    \left< F \right>&=v \\
    \left< 1-F \right>& = 1-v.
\end{align}
We note the following properties of the random function: 

Now we apply the machinery to our system of equations. We declare phase 1 PuO and phase 2 the homogeneously averaged combustible mix. We declare all functions of $\mathbf{r}$ as random functions. 

\subsection{Stochastic Heat Conduction}
Equation \ref{SCE} becomes
\begin{align}
    \rho_1 C_{p1} \frac{\partial T_1}{\partial t}F +\rho_2 C_{p2} \frac{\partial T_2}{\partial t}(1-F)=&-F \mathbf{\nabla} \cdot \mathbf{J}_1 - \mathbf{J}_1 \cdot \mathbf{\nabla}F -(1-F) \mathbf{\nabla} \cdot \mathbf{J}_2 +\mathbf{J}_2 \cdot \mathbf{\nabla}F \notag \\ &+S_1 F + S_2 (1-F)
\end{align}
where we have introduced $\mathbf{J}(\mathbf{r})=-K(\mathbf{r})\mathbf{\nabla}T(\mathbf{r},t)$. Multiplying by $F$ and $1-F$ we get, respectively 
\begin{align}
    \rho_1 C_{p1}\frac{\partial T_1}{\partial t}F&=-F\nabla \cdot \mathbf{J}_1-\mathbf{J}_1 \cdot F\nabla F+\mathbf{J}_2 \cdot F\nabla F + S_1 F \\
    \rho_2 C_{p2}\frac{\partial T_2}{\partial t}(1-F)&=-\mathbf{J}_1 \cdot (1-F)\nabla F-(1-F)\nabla \cdot \mathbf{J}_2+\mathbf{J}_2 \cdot (1-F)\nabla F + S_2 (1-F).
\end{align}
To deal with the above, we perform a statistical average. Clearly we must take care to deal with $\left< F\nabla F \right>$. From first principles 
\begin{align}
    \left< F(x) F'(x)\right>&=\lim _{h\rightarrow 0}\left< F(x) \frac{1}{h} \left(F(x)-F(x-h)\right)\right> \notag \\
    &=\lim_{h\rightarrow 0} \frac{1}{h}\left( \left< F(x)\right>-\left< F(x)F(x-h)\right> \right).
\end{align}
To deal with $F(x)F(x-h)$, we introduce the \textbf{correlation function}, defined as 
\begin{equation}
    R(h)=\frac{\left< F(x)F(x-h)\right>-\left< F(x)\right>^2}{\left< F(x)^2\right>-\left< F(x)\right>^2}=\frac{\left< F(x)F(x-h)\right>-v^2}{v-v^2}.
\end{equation}
Hence 
\begin{align}
    \left< F(x)F'(x)\right>&=\lim_{h\rightarrow 0} \frac{1}{h}\left(v-v(1-v)R(h)-v^2\right) \notag \\
    &=\lim_{h\rightarrow 0} \frac{1}{h} \left( v-v^2-v(1-v)\left(R(0)+hR'(0)+o(h^2)\right)\right) \notag \\ 
    &=-R'(0)v(1-v) +\lim_{h\rightarrow 0}\left( v-v^2-R(0)v +R(0)v^2\right).
\end{align}
Evidently we must impose $R(0)=1$, so the limit converges, and thus we have 
\begin{equation}
    \left< F(x)F'(x)\right>=-R'(0)v(1-v).
\end{equation}
This generalises trivially to give 
\begin{equation}
    \left< F(\mathbf{r})\nabla F(\mathbf{r})\right>=-v(1-v) \nabla R(\mathbf{r}) \Bigr|_{\mathbf{r=0}}.
\end{equation}
The other problem of note is the source term. $S_{fission}$ is trivial, but 
\begin{equation}
S_{combustion}=q(\mathbf{r})k(\mathbf{r})\omega(\mathbf{r},t)\exp{\left(-\frac{E_{act}(\mathbf{r})}{R_g T(\mathbf{r},t)}\right)}
\end{equation}
poses a challenge. After multiplying by $F$ and $1-F$, and statistical averaging,
\begin{align}
    vS_{combustion,1}&=q_1 k_1 \omega_1 \left<F \exp{\left(-\frac{E_{act}(\mathbf{r})}{R_g T(\mathbf{r},t)}\right)}\right> \notag \\
    (1-v)S_{combustion,2}&=q_2 k_2 \omega_2 \left<(1-F) \exp{\left(-\frac{E_{act}(\mathbf{r})}{R_g T(\mathbf{r},t)}\right)}\right>.
      \end{align}
To deal with this, we must first get the exponent in a more tractable form. We have 
\begin{align}
    \frac{E_{act}(\mathbf{r})}{T(\mathbf{r},t)}&=\frac{E_{a1}F+E_{a2}(1-F)}{T_1F+T_2(1-F)} = \frac{E_{a2}+\Delta E_a F}{T_2+\Delta TF} \notag \\
    &=\frac{E_{a2}}{T_2+\Delta TF}+\frac{\Delta E_a}{\Delta T} \frac{F}{\frac{T_2}{\Delta T}+F} \notag \\
    &=\frac{E_{a2}}{T_2+\Delta TF} + \frac{\Delta E_a}{\Delta T}- \frac{T_2}{\Delta T} \frac{\Delta E_a}{\Delta T} \frac{1}{\frac{T_2}{\Delta T}+F} \notag \\
    &=\frac{\Delta E_a}{\Delta T}+\left( E_{a2}-\frac{\Delta E_a}{\Delta T}T_2\right) \frac{1}{T_2 + \Delta T F} \notag \\
    &=\frac{\Delta E_a}{\Delta T}+\left( E_{a2}-\frac{\Delta E_a}{\Delta T}T_2\right) \int _0 ^\infty du e^{-u(T_2+\Delta TF)}.
\end{align}
We note that $E_a/R_g=T_a$. Now, exponentiating: 
\begin{align}
    \exp \left(- \frac{T_a}{T}\right)=e^{-\frac{\Delta T_a}{\Delta T}}\exp{\left[ -\left(T_{a2}-\frac{\Delta T_a}{\Delta T}T_2 \right)\int_0 ^{\infty}due^{-u(T_2+\Delta TF)} \right]}
\end{align}
And expanding in a power series: 
\begin{align}
     \exp \left( -\frac{T_a}{T}\right)&=e^{-\frac{\Delta T_a}{\Delta T}} \sum_{n=0}^\infty \frac{(-1)^n}{n!}  \left(T_{a2}-\frac{\Delta T_a}{\Delta T}T_2 \right)^n \left( \int_0 ^{\infty}due^{-u(T_2+\Delta TF)}\right)^n \notag \\
     &=e^{-\frac{\Delta T_a}{\Delta T}} \sum_{n=0}^\infty \frac{(-1)^n}{n!}  \left(T_{a2}-\frac{\Delta T_a}{\Delta T}T_2 \right)^n \left[ \int_0 ^{\infty}due^{-uT_2}\left(1+\sum_{m=1}^\infty \frac{(-1)^m}{m!}(u\Delta T)^m F^m\right)\right]^n \notag \\
     &=e^{-\frac{\Delta T_a}{\Delta T}} \sum_{n=0}^\infty \frac{(-1)^n}{n!}  \left(T_{a2}-\frac{\Delta T_a}{\Delta T}T_2 \right)^n \left[ \int_0 ^{\infty}due^{-uT_2}\left(1+F\sum_{m=1}^\infty \frac{(-1)^m}{m!}(u\Delta T)^m \right)\right]^n
\end{align}
In the above we used the slightly modified $e^{F}=1+\sum_{m=1}^\infty \frac{F^m}{m!}$. Making use of $F^n=F$ and $(1-F)^n=1-F$
\begin{align}
    F\exp{\left( -\frac{T_a}{T}\right)}&=e^{-\frac{\Delta T_a}{\Delta T}} \sum_{n=0}^\infty \frac{(-1)^n}{n!}  \left(T_{a2}-\frac{\Delta T_a}{\Delta T}T_2 \right)^n \left[ \int_0 ^{\infty}due^{-uT_2}\left(F+F\sum_{m=1}^\infty \frac{(-1)^m}{m!}(u\Delta T)^m \right)\right]^n \notag \\
    &=Fe^{-\frac{\Delta T_a}{\Delta T}} \sum_{n=0}^\infty \frac{(-1)^n}{n!}  \left(T_{a2}-\frac{\Delta T_a}{\Delta T}T_2 \right)^n \left[ \int_0 ^{\infty}due^{-uT_2}\left(1+\sum_{m=1}^\infty \frac{(-1)^m}{m!}(u\Delta T)^m \right)\right]^n \notag \\
    &=Fe^{-\frac{\Delta T_a}{\Delta T}}\exp{\left[ -\left( T_{a2}-\frac{\Delta T_a}{\Delta T}T_2\right) \int_0 ^\infty du e^{ -u(T_2+\Delta T)} \right]} \notag \\
    &=Fe^{-\frac{\Delta T_a}{\Delta T}}\exp{\left[ -\left( T_{a2}-\frac{\Delta T_a}{\Delta T}T_2\right) \frac{1}{T_2 + \Delta T} \right]}=Fe^{-\frac{T_{a1}}{T_1}}
\end{align}
and similarly 
\begin{equation}
    (1-F)\exp{\left( -\frac{T_a}{T}\right)}=(1-F)e^{-\frac{T_{a2}}{T_2}}.
\end{equation}
Hence we have 
\begin{align}
    vS_{combustion,1}&=q_1 k_1 \omega_1ve^{-\frac{T_{a1}}{T_1}}  \\
    (1-v)S_{combustion,2}&=q_2 k_2 \omega_2(1-v)e^{-\frac{T_{a2}}{T_2}}.
\end{align}
So our stochastic pyrolysis equations amount to 
\begin{align}
    \rho_1 C_{p1} \frac{\partial T_1}{\partial t}&= -\nabla \cdot \mathbf{J}_1 +(1-v)(\mathbf{J}_1-\mathbf{J}_2)\cdot \nabla R(\mathbf{r}) |_{\mathbf{r}=\mathbf{0}} + E_f \int_0 ^ \infty dE \Sigma_f \phi_1 + q_1 k_1 \omega_1 e^{-\frac{E_{a1}}{R_g T_1}} \\
    \rho_2 C_{p2} \frac{\partial T_2}{\partial t}&= -\nabla \cdot \mathbf{J}_2 -v(\mathbf{J}_1-\mathbf{J}_2)\cdot \nabla R(\mathbf{r}) |_{\mathbf{r}=\mathbf{0}} + E_f \int_0 ^ \infty dE \Sigma_f \phi_2 + q_2 k_2 \omega_2 e^{-\frac{E_{a2}}{R_g T_2}}
\end{align}
with the currents given by 
\begin{align}
    \mathbf{J}_1 &= -K_1 \nabla T_1 +K_1 (1-v)(T_1 -T_2) \nabla R(\mathbf{r})|_{\mathbf{r}=\mathbf{0}} \\
    \mathbf{J}_2 &= -K_2 \nabla T_2 -K_2 v(T_1 -T_2) \nabla R(\mathbf{r})|_{\mathbf{r}=\mathbf{0}}.
\end{align}

\subsection{Stochastic Neutron Diffusion}
There are no new concerns in the NDE that have not been dealt with in the previous section. The stochastic NDE takes the form
\begin{align}
    \frac{1}{\upsilon(E)} \frac{\partial \phi_1 (\bold{r},E,t)}{\partial t}&= -\nabla \cdot \mathbf{Q}_1+(1-v)(\mathbf{Q}_1-\mathbf{Q}_2)\cdot \nabla R(\mathbf{r})|_{\mathbf{r}=\mathbf{0}}-\Sigma_{a1} (\bold{r},E)\phi_1 \notag \\ &+\int_{0}^{\infty}\Sigma_{s1}(\bold{r}, E'\rightarrow E)\phi_1(\bold{r},E',t)dE'  +\chi_p(E) \int_0 ^{\infty} \nu(E') (1-\beta)\Sigma_{f1} (\bold{r},E')\phi_1(\bold{r},E',t)dE' \notag \\ 
    &+\sum_{i=1}^{I} \chi_{d,i}(E)\lambda_i C_{i1}(\bold{r},t) 
    \\ \notag \\ 
    \frac{1}{\upsilon(E)} \frac{\partial \phi_2 (\bold{r},E,t)}{\partial t}&= -\nabla \cdot \mathbf{Q}_2-v(\mathbf{Q}_1-\mathbf{Q}_2)\cdot \nabla R(\mathbf{r})|_{\mathbf{r}=\mathbf{0}}-\Sigma_{a2} (\bold{r},E)\phi_2 \notag \\ &+\int_{0}^{\infty}\Sigma_{s2}(\bold{r}, E'\rightarrow E)\phi_2(\bold{r},E',t)dE'  +\chi_p(E) \int_0 ^{\infty} \nu(E') (1-\beta)\Sigma_{f2} (\bold{r},E')\phi_2(\bold{r},E',t)dE' \notag \\ 
    &+\sum_{i=1}^{I} \chi_{d,i}(E)\lambda_i C_{i2}(\bold{r},t)
\end{align}
with precursor equations 
\begin{align}
    \frac{\partial C_{i1}({\bold{r},t)}}{\partial t}&=\beta_i \int_0^\infty \nu(E) \Sigma_{f1} (\bold{r},E) \phi_1(\bold{r}, E,t)dE-\lambda_i C_{i1} \\
    \frac{\partial C_{i2}({\bold{r},t)}}{\partial t}&=\beta_i \int_0^\infty \nu(E) \Sigma_{f2} (\bold{r},E) \phi_2(\bold{r}, E,t)dE-\lambda_i C_{i2}
\end{align}
and currents 
\begin{align}
    \mathbf{Q}_1 &= -D_1 \nabla \phi_1 +D_1 (1-v)(\phi_1 -\phi_2) \nabla R(\mathbf{r})|_{\mathbf{r}=\mathbf{0}} \label{Q1}\\
    \mathbf{Q}_2 &= -D_2 \nabla \phi_2 -D_2 v(\phi_1 -\phi_2) \nabla R(\mathbf{r})|_{\mathbf{r}=\mathbf{0}}. \label{Q2}
\end{align}
We justify stochastically distributing the NDE as large chunks of fissile material lumped together is where the risk of criticality is highest.


\section{Approximations, Assumptions}
\subsection{The Correlation Function}
The only restraint on $R(\mathbf{r})$ is that $R(\mathbf{0})=1$, so we have a lot of freedom in choosing the correlation function. We choose to let our intuition inform us of a physically reasonable form. A simple choice is 
\begin{equation}
    R(\mathbf{r})=e^{-\boldsymbol{\mu} \cdot \mathbf{r}}.
\end{equation}

\subsection{Thermal Approximations}

\subsection{Diffusion Approximations}
The first approximation to implement is the \textbf{multigroup approximation}, where we split up integrals over energy as 
\begin{equation}
    \int_{E_g}^{E_{g-1}}f(E)dE=f_g .
\end{equation}
Equaion \eqref{TDNDE} becomes 
\begin{align}
    \frac{1}{\upsilon_g}\frac{\partial \phi_g(\mathbf{r},t)}{\partial t}=& \nabla \cdot (D_g(\mathbf{r})\nabla \phi_g)-\Sigma^g_a(\mathbf{r})\phi_g+\sum_{g'=1}^G\Sigma_s^{g'\rightarrow g}(\mathbf{r}) \phi_{g'}(\mathbf{r},t) \\& +\chi_p ^g \sum_{g'=1}^G \nu_{g'} (1-\beta)\Sigma_f^{g'}(\mathbf{r})\phi_{g'}(\mathbf{r},t)+\sum_{i=1}^I \chi_{d,i}^g \lambda_i C_i (\mathbf{r},t)
\end{align}
with precursors
\begin{equation}
    \frac{\partial C_i(\mathbf{r},t)}{\partial t}=\beta_i \sum_{g=1}^G \nu_g \Sigma_f^g(\mathbf{r})\phi_g(\mathbf{r},t)- \lambda_i C_i.
\end{equation}
Let us now introduce some cleaner notation. We define the \textbf{fission operator} $F$ and the \textbf{absorption operator} $A$ such that 
\begin{align}
    F\phi=&\sum_{g'=1}^G \nu_{g'} \Sigma_f^{g'}(\mathbf{r})\phi_{g'}(\mathbf{r},t) \\
    (A\phi)_g =&\nabla \cdot (D_g(\mathbf{r})\nabla \phi_g)-\Sigma^g_a(\mathbf{r})\phi_g+\sum_{g'=1}^G\Sigma_s^{g'\rightarrow g}(\mathbf{r}) \phi_{g'}(\mathbf{r},t) 
\end{align}
where $\phi$ without an index is a column vector comprised of $\phi_g$. Hence 
\begin{align}
    \frac{1}{\upsilon} \frac{\partial \phi}{\partial t}=&A\phi+(1-\beta)\chi_p F\phi+\sum_{i=1}^{I}\lambda_i \chi_{d,i}C_i \\
    \frac{\partial C_i}{\partial t} =&\beta_i F \phi -\lambda_i C_i.
\end{align}


When comparing with the stochastic diffusion equations we found earlier, we see the need to define a \textbf{stochastic loss operator} $L_i$ for each flux field $\phi_i$. (We note that these flux fields are still column vectors comprised of $G$ energy groups). Let us write out the equation for $\phi_1$:
\begin{equation}
    \frac{1}{\upsilon} \frac{\partial \phi_1}{\partial t}= -\nabla \cdot \mathbf{Q}_1 +(1-v)(\mathbf{Q}_1-\mathbf{Q}_2)\cdot \nabla R|_{\mathbf{r}=\mathbf{0}}+A_1\phi_1 +(1-\beta)\chi_p F_1\phi_1 + \sum_{i=1}^{I}\lambda_i \chi_{d,i}C_{i1}
\end{equation}
Note that the absorption operator $A$ no longer contains the derivative term. Let us now define
\begin{equation}
    L_1 \phi_1=-\nabla \cdot \mathbf{Q}_1 -v_2 (\mathbf{Q}_2-\mathbf{Q}_1)\cdot \nabla R|_{\mathbf{r}=\mathbf{0}}.
\end{equation} \\ 

We note that we can generalise to an arbitrary number of phases as follows. The current for each phase becomes:
\begin{equation}
    \mathbf{Q}_i=-D_i \nabla \phi_i +D_i  \sum_j v_j (\phi_i - \phi_j)\nabla R_{ij}(\mathbf{r})|_{\mathbf{r}=0}.
\end{equation}
So that we can define 
\begin{equation}
     L_i \phi_i=-\nabla \cdot \mathbf{Q}_i - \sum_k v_k (\mathbf{Q}_k-\mathbf{Q}_i) \nabla R_{ik}|_{\mathbf{r}=\mathbf{0}}
\end{equation}
and thus our equations are  
\begin{align}
    \frac{1}{\upsilon} \frac{\partial \phi_i}{\partial t} &=L_i \phi_i +A_i\phi_i +(1-\beta)\chi_p F_i\phi_i+\sum_{j=1}^I\lambda_j \chi_{d,j}C_{ji} \\
    \frac{\partial C_{ji}}{\partial t}&=\beta_j F_i\phi_i-\lambda_j C_{ji}.
\end{align}

We also adopt an adiabatic (or quasistatic) approach, in which we write the flux as $\phi_g(\mathbf{r},t)=P(t)\psi_g(\mathbf{r},t)$, the precursor concentration as $C_{ji}(\mathbf{r},t)=P(t)\hat{C}_{ji}(\mathbf{r},t) $ and assume $\frac{\partial \psi}{\partial t}\approx 0$, $\frac{\partial \hat{C}_{ji}}{\partial t}\approx 0$ (we justify this as the timescale on which the shape changes due to thermal effects is much greater than the neutron lifetime). We will advance in time with the point kinetics equations, and use a static flux solve at each timestep to update the parameters. We normalise the flux 
\begin{equation}
    \sum_i \sum_g \int _V \psi_{ig}(\mathbf{r})d\mathbf{r}=1
\end{equation}
so that $P(t)$ is the total neutron population. We can make a further approximation to the precursor shape function $\hat{C_{ji}}(\mathbf{r})$
\begin{equation}
    \hat{C_{ji}}=\frac{\beta_j}{\lambda_j}F_i \psi_i.
\end{equation}
Thus we have for the precursor source term
\begin{equation}
    \sum_j \lambda_j \chi_{d,j} \hat{C}_{ji}=\sum_j \beta_j \chi_{d,j}
\end{equation}
and we can write 
\begin{equation}
    \chi_{eff}=(1-\beta)\chi_p +\sum_j \beta_j \chi_{d,j}.
\end{equation}
The shape functions satisfy the static (stochastic) eigenvalue problem at each timestep: 
\begin{equation}
    (L_i +A_i) \psi_i +\frac{1}{k} \chi_{eff} F_i \psi_i=0.
\end{equation}
This will be solved with updated cross sections from temperature fields and volume fractions. The eigenvalue $k$ extracted gives the reactivity as 
\begin{equation}
    \rho(t)= \frac{k(t)-1}{k(t)}.
\end{equation}
The prompt neutron lifetime is computed as 
\begin{equation}
    \Lambda (t)=\frac{\sum_i \sum_g \int_V \psi_{ig}/v_g d\mathbf{r}}{\sum_i \sum_g \nu_g \Sigma_{fi}^g\int_V \psi_{ig} d\mathbf{r}}.
\end{equation}
And finally the amplitude is evolved in time via the point kinetics equations 
\begin{align}
    \Lambda (t)\frac{dP}{dt}&=(\rho(t)-\beta)P(t)+\sum_{j=1}^I\lambda_j \hat{C}_j (t) \\ 
    \frac{d\hat{C}_j }{dt}&=\frac{\beta_j}{\Lambda (t)} P(t) -\lambda_j \hat{C}_j(t).
\end{align}





% Substituting in and rearranging we get 
% \begin{align}
%     \frac{1}{\upsilon} \frac{\dot{P}}{P}\psi_i &=L_i \psi_i +A_i\psi_i +(1-\beta)\chi_p F_i\psi_i+\sum_{j=1}^I\lambda_j \chi_{d,j}\hat{C}_{ji} \\
%     \frac{\dot{P}}{P}\hat{C}_{ji}&=\beta_j F_i\psi_i-\lambda_j \hat{C}_{ji}.
% \end{align}
% We will now take the inner product with the adjoint flux $\psi^\dagger_j (\mathbf{r})$ over space and energy groups in order to extract scalar solutions. We must also sum over all phases:
% \begin{align}
%     \dot{P}\sum_i \left<\psi_i^\dagger , \psi _i \right> 
% \end{align}
    
\end{document}


